
This book grew out of a problem I faced when writing
my PhD thesis on flutter of windmill blades.
I needed to study the effects of varying properties such
as blade stiffness, mass properties, and aerodynamic properties,
but existing solution techniques made studies like these
very tedious.

A classmate suggested two techniques that seemed like a good
fit for my problem: continuation methods and automatic differentiation.
The key advantage continuation methods have over traditional
eigenvalue methods is the use of derivatives of the equations.
With these derivatives continuation methods can trace continuous curves,
in addition to solving flutter problems that are difficult or
impossible with traditional techniques.

When I joined Boeing after graduation I showed some results to
a senior flutter engineer who convinced me to let others
try the program I developed for my thesis.
It was an immediate success, largely due to its ability to
perform parameter studies.
During my 32 year career at Boeing I added numerous features,
but some of the more important ones had to wait until I retired.

In retirement I wrote a completely new implementation of the method in
modern C++, designed around nonlinear flutter, extensible plugins,
automatic differentiation, and parallel execution.
This program, called \emph{Fina}, serves as the computational foundation
for this book and is available as open-source software.
To generate the example problems, I added a small beam modeling capability,
which is ideal for demonstrating the solver's capabilities and providing
reproducible test cases for the journal articles I have written.

I have chosen to publish this book as open-access.
My ultimate goal is simply to share this knowledge and improve the
discipline of flutter analysis -- not to become fabulously wealthy.
